\chapter{Normalización}

\section{¿Por qué normalizar?}

Normalización elimina redundancia y anomalías de actualización. Una tabla desnormalizada puede tener: anomalía de inserción (no puedo insertar sin datos extras), de actualización (cambio en un lugar no se propaga) y de borrado (pierdo datos al borrar fila).

\section{Formas normales}

\begin{itemize}
  \item \textbf{1NF:} valores atómicos (sin listas/arrays en celdas), sin filas duplicadas, PK definida
  \item \textbf{2NF:} en 1NF + sin dependencias parciales (atributos dependen de toda la PK, no solo parte)
  \item \textbf{3NF:} en 2NF + sin dependencias transitivas (atributo no-clave no depende de otro atributo no-clave)
  \item \textbf{BCNF:} en 3NF + todo determinante es superkey (más estricta que 3NF)
\end{itemize}

\section{Dependencias funcionales}

$X \rightarrow Y$: conocer X determina Y. Si PK = \{A, B\} y $A \rightarrow C$ (C no depende de B), hay dependencia parcial → viola 2NF.

\begin{lstlisting}[caption={Tabla que viola 2NF y su corrección}]
-- VIOLA 2NF: PK=(id_pedido, id_producto) pero nombre_producto
--            depende solo de id_producto (dependencia parcial)
CREATE TABLE pedido_items_malo (
    id_pedido       INTEGER,
    id_producto     INTEGER,
    nombre_producto VARCHAR(100),
    cantidad        INTEGER,
    PRIMARY KEY (id_pedido, id_producto)
);

-- CORRECTO: separar producto a su propia tabla
CREATE TABLE productos (
    id     SERIAL PRIMARY KEY,
    nombre VARCHAR(100) NOT NULL
);

CREATE TABLE pedido_items (
    id_pedido   INTEGER,
    id_producto INTEGER REFERENCES productos(id),
    cantidad    INTEGER NOT NULL,
    PRIMARY KEY (id_pedido, id_producto)
);
\end{lstlisting}

\begin{lstlisting}[caption={Tabla que viola 3NF y su corrección}]
-- VIOLA 3NF: id_depto -> nombre_depto (dependencia transitiva)
CREATE TABLE empleados_malo (
    id_empleado  SERIAL PRIMARY KEY,
    nombre       VARCHAR(100),
    id_depto     INTEGER,
    nombre_depto VARCHAR(80)
);

-- CORRECTO: tabla separada para departamentos
CREATE TABLE departamentos (
    id     SERIAL PRIMARY KEY,
    nombre VARCHAR(80) NOT NULL
);

CREATE TABLE empleados (
    id       SERIAL PRIMARY KEY,
    nombre   VARCHAR(100),
    id_depto INTEGER REFERENCES departamentos(id)
);
\end{lstlisting}

\begin{entrevistabox}
\begin{enumerate}
  \item ¿Cuál es la diferencia entre 2NF y 3NF?
  \item ¿Cuándo desnormalizarías intencionalmente una tabla?
  \item Identifica la forma normal: \texttt{(orden\_id, producto\_id, categoria\_producto, cantidad)} con PK=(orden\_id, producto\_id).
\end{enumerate}
\end{entrevistabox}
